% chapters/ch20_curvature.tex
\chapter{Triplets as Ultrametric Curvature Observations}
\label{ch:curvature}

\epigraph{%
  To measure a tree, you need not see the whole tree.
  You need only see which paths branch.%
}{\textit{---Flyxion}}

This chapter develops the geometric interpretation of the splitting oracle. We showed in Chapter~\ref{ch:metrics} that every hierarchical tree corresponds to an ultrametric, and in Chapter~\ref{ch:triplets} that every triplet relation corresponds to a specific ultrametric inequality. Combining these two observations gives a clean geometric interpretation: each oracle query is a local measurement of the curvature structure of the ultrametric space defined by $T^*$. The oracle does not reveal the actual distance values---only the ordering of distances for each triple---and this minimal information is precisely what is needed to determine the tree topology, as Theorem~\ref{thm:dense_triplets} established. The analogy with curvature is not merely metaphorical: in negatively curved spaces, geodesics diverge in a tree-like pattern, and the branching structure of that pattern is what the oracle measures.

%% ── Figure: geodesic branching and triplet observation ───────
\begin{figure}[ht]
\centering
\begin{tikzpicture}[scale=1.05]
% Y-shaped tree as geodesic merging
\coordinate (center) at (0,0);
\coordinate (u) at (-1.8, 1.5);
\coordinate (v) at (1.8, 1.5);
\coordinate (w) at (0, -2.0);

\draw[nodeblue!70, very thick] (u) -- (center);
\draw[nodeblue!70, very thick] (v) -- (center);
\draw[nodegreen!70, very thick] (center) -- (w);

% Dots at endpoints
\node[leafnode, label=above:{$u$}] at (u) {};
\node[leafnode, label=above:{$v$}] at (v) {};
\node[leafnode, label=below:{$w$}] at (w) {};
\node[internalnode, label=right:{\footnotesize branching point}] at (center) {};

% Distance annotations
\node[font=\footnotesize, nodeblue!80] at (-0.9, 0.85)
  {$d(u,v)$ short};
\node[font=\footnotesize, nodegreen!80] at (0.8, -1.0)
  {$d(u,w)=d(v,w)$ long};

% Oracle output annotation
\draw[->, >=stealth, nodered!70, thick] (2.5, 0) -- (3.5, 0);
\node[font=\small, nodered!80, right] at (3.5, 0.2)
  {Oracle returns: $w$};
\node[font=\footnotesize, nodegray, right] at (3.5,-0.35)
  {(i.e., $(u,v)\mid w$)};
\node[font=\footnotesize, nodegray, right] at (3.5,-0.75)
  {$d(u,v) < d(u,w) = d(v,w)$};
\end{tikzpicture}
\caption{The oracle query on $(u,v,w)$ reveals the branching geometry of the ultrametric space: it identifies which of the three pairs shares the lower branching point (shorter ultrametric distance). The Y-shaped tree makes this geometric: the oracle identifies the pair $\{u,v\}$ whose geodesics share the most path before diverging.}
\label{fig:curvature}
\end{figure}

\section{The Oracle as a Geometric Instrument}

\begin{theorem}[Oracle as curvature measurement]
The splitting oracle, queried on $(u,v,w)$, reveals which of the three strict ultrametric inequality patterns holds:
\begin{align*}
(u,v) \mid w &\iff d_{T^*}(u,v) < d_{T^*}(u,w) = d_{T^*}(v,w),\\
(u,w) \mid v &\iff d_{T^*}(u,w) < d_{T^*}(u,v) = d_{T^*}(v,w),\\
(v,w) \mid u &\iff d_{T^*}(v,w) < d_{T^*}(u,v) = d_{T^*}(u,w).
\end{align*}
\end{theorem}

\begin{proof}
Immediate from Proposition~\ref{prop:ultrametric_lca} and the isosceles property (Proposition~\ref{prop:isosceles}).
\end{proof}

\section{Sparse Curvature Observations}

\begin{corollary}
The splitting oracle provides exactly the information needed to reconstruct the ultrametric tree topology, and no more: it reveals the branching structure but not the merge heights.
\end{corollary}

Knowing all triplet relations for $n$ leaves is equivalent to knowing the tree topology but not the height function.

\section{Exercises}

\begin{enumerate}
  \item Describe the oracle's output in terms of the ultrametric distance matrix for a specific tree on five leaves.
  \item Show that knowing distance orderings (not values) suffices for topology reconstruction.
  \item Discuss in what sense the oracle measures ``discrete curvature'' and how this relates to curvature in Riemannian geometry.
\end{enumerate}
